%document class
\documentclass[10pt,oneside]{book}
%%%% Page Info + Commands %%%%%{

%packages
\usepackage{pgfplots}
\pgfplotsset{compat=1.18}
\usepackage{amsmath}
\usepackage{amsfonts}
\usepackage{charter}
\usepackage{amsthm,letltxmacro}
\usepackage{amssymb}
\usepackage{fancyhdr}
\usepackage{float}
\usepackage[most]{tcolorbox}
\usepackage{enumerate}
\usepackage{enumitem}
\usepackage[dvipsnames]{xcolor}
\usepackage{changepage}
\usepackage{mdframed}
\usepackage{graphicx}

% SMALL PAGE COMMAND
\newcommand{\smallpage}{  \setlength \oddsidemargin{.25in}
            \setlength \textwidth{5.5in}
            \setlength \topmargin{0in}
            \setlength \textheight{9in}}


% Commands for sets of numbers
\newcommand{\Z}{\mathbb{Z}}\newcommand{\R}{\mathbb{R}}\newcommand{\N}{\mathbb{N}}\newcommand{\Q}{\mathbb{Q}}\newcommand{\C}{\mathbb{C}}
\newcommand{\real}[1]{\text{Re}\left(#1\right)}
\newcommand{\im}[1]{\text{Im}\left(#1\right)}

% gordie spacing and boxing commands
\newcommand{\gspc}{\vspace{0.13cm}} % 0.5 space
\newcommand{\gline}[0]{\gspc\\} % 1.5 space
\newcommand{\gbline}{\gspc\gspc\\} % double space
\newcommand{\gbox}[1]{\fbox{\parbox{\linewidth}{#1}}} % multiline box command
\newcommand{\gmod}[1]{\,(\text{mod}\,#1)}


\newcommand{\gimage}[3]{
\begin{figure}[H]
    \centering
    \includegraphics[width=#2]{#1}
    \caption{#3}
    \label{fig:#1}
\end{figure}
}

\newcommand{\centerit}[1]{{\begin{center} #1 \end{center}}}
\newcommand{\ul}[1]{\underline{#1}}

%%%%%%%% command for graphics %%%%%%%%%%%%%

%}

% COLOR BOX
\newmdenv[backgroundcolor=gray!8,
  linewidth=0pt, innerleftmargin=0pt, innerrightmargin=0pt,
  skipabove=\baselineskip, skipbelow=\baselineskip
]{coloredadjust}

% PROOF
\LetLtxMacro\oldproof\proof
\let\endoldproof\endproof
\renewenvironment{proof}[1][\proofname]
  {\vspace{0.2cm}\begin{adjustwidth}{2em}{2em}
   \oldproof[#1]}
  {\endoldproof
\end{adjustwidth}}

% PROBLEM
\newenvironment{problem}[1]
  {\vspace{-0.1cm}\begin{coloredadjust}\begin{adjustwidth}{2em}{2em}
   \textit{Problem. }#1}
  {\endoldproof
\end{adjustwidth}\end{coloredadjust}\gspc}

%% DEFINITION
\newtcbtheorem[]{definition}{Definition}{enhanced,
  colback=blue!2, colframe=blue!50!black, coltitle=blue!70!black,
  fonttitle=\bfseries,
  boxrule=0pt, toprule=2pt, bottomrule=1pt, leftrule=1pt, rightrule=1pt, arc=1pt, outer arc=1pt,
  attach title to upper, title={#1},
}{dfn}

%% CRITERION
\newtcbtheorem[]{criterion}{Criterion}{enhanced,
  colback=magenta!3, colframe=magenta!60!black, coltitle=magenta!20!black,
  fonttitle=\bfseries,
  boxrule=4pt, toprule=2pt, bottomrule=1pt, leftrule=1pt, rightrule=1pt, arc=1pt, outer arc=1pt,
  attach title to upper, title={#1},
}{ctn}

%% COROLLARY
\newtcbtheorem[]{corollary}{Corollary}{enhanced,
  colback=orange!3, colframe=orange!80!black, coltitle=orange!50!black,
  fonttitle=\bfseries,
  boxrule=4pt, toprule=1pt, bottomrule=1pt, leftrule=1pt, rightrule=1pt, arc=5pt, outer arc=5pt,
  attach title to upper, title={#1},
}{ctn}

%% THEOREM
\newtcbtheorem[]{theorem}{Theorem}{enhanced,
  colback=green!2, colframe=green!40!black, coltitle=green!20!black,
  fonttitle=\bfseries,
  boxrule=4pt, toprule=3pt, bottomrule=1pt, leftrule=1pt, rightrule=1pt,arc=5pt, outer arc=5pt,
  attach title to upper, title={#1},
}{thm}

%%%% Page Setup %%%%%%%
\smallpage\sloppy
\pagestyle{fancy}
\lhead{\large{\textbf{Feb 6th{}-{}Notes}}}
%\rfoot{\thepage/\pageref{LastPage} }
\setlength{\headheight}{10pt}

\begin{document}

\newcommand{\cg}[1]{\color{OliveGreen}#1\color{white}}
\newcommand{\cb}[1]{\color{BlueViolet}#1\color{white}}

\subsection*{Introduction}
First, we started with introductions for Professor Stoertz. Professor Stoertz has two cats: Reese and Hank. Also two dogs: a husky and a part-husky named Lulu.
\subsection*{Class Expectations}
There's a couple of expectations student should be aware of:
\begin{itemize}
  \item AI Use is strongly discouraged, but if used, students should cite it. 
  \item Email professor stoertz frequently. 
  \item Now for some learning goals: \begin{itemize}
    \item Adopt a growth mindset.
    \item Work to improve teamwork skills.
    \item Leard the methods and language of Complex Analysis. 
  \end{itemize}
  \item Components of the grade: \begin{itemize}
    \item \textbf{10\% Active Learning and Engagemenet}
    \item \textbf{10\% Daily Preview Activities}
    \item \textbf{35\% Weekly written Homework assignments} \begin{itemize}
      \item Assignments will be posted onOverleaf and will be due on Fridays (\textit{mostly}) by 8:00pm
      \item Then, we grade ourselves on the following criteria: \begin{itemize}
        \item Accuracy of grade you have given yourself
        \item Accuracy and quality of your grading
        \item A good-faith attempt at every problemClear communication of solutions
        \item A collaboration statement
        \item Reflections
        \item *\textit{Lowest homework grade will be dropped}
      \end{itemize}
    \end{itemize}
    \item \textbf{HW \#0} - If we drop in by Feb. 20th, we get two more dropped preview assignments and more more dropped homework
    \item \textbf{30\% Three Midtern Exams} - A portion of these exams could be take-home.
    \item \textbf{15\% Final Project} - Three or four student teams on a project with a 10 min presentation during the final period. 
  \end{itemize}
  \item Also remember to have fun and learn. 
\end{itemize}
\subsection*{1.1 - The Algebra of Complex Numbers}
First, we cover the basic sets of numbers
\begin{align*}
  \N &= \{0, 1, 2\}\\
  \Z &= \{-2, -1, 0, 1, 2\}\\
  \Q &= \{\frac{a}b \mid a\in \Z, b\in \N \backslash\{0\}\}\\
  \R &= \text{ Solving Equations like } x^2 - 2 = 0\\
  \C &= \{a+bi\mid a,b\in\R\}
\end{align*}
\newpage
\begin{definition}{Complex Numbers}{}\gline{}
  Denote by $i$ the number $\sqrt{-1}$. We define the \ul{set of complex numbers}
  \begin{align*}
    \C =\{a+bi\mid a,b\in\R\}
  \end{align*}
  \centerit{\ul{\textbf{Equality}}}
  Two complex numbers $a+bi$ and $c+di$ are said to be equal if and only if:
  \begin{align*}
    a = c \text{ and } b = d
  \end{align*}
  \centerit{\ul{\textbf{Addition}}}
  We add complex numbers with the following formula;
  \begin{gather*}
    (a+bi) + (c + di) = (a+c) + (c + d)i
  \end{gather*}
  \centerit{\ul{\textbf{Multiplication}}}
  We multiply complex numbers with the following formula;
  \begin{align*}
    (a+bi)(c+di)&= (ac + adi + bci + bdi^2)\\
    &= (ac-bd) + (ad+bc)i
  \end{align*}
\end{definition}
We will cover all the axioms of $\C$ at a later point. 
\begin{theorem}{$\C$ is a field}{}\gline{}
  \textit{Note: }$\C$ is not an ordered field, meaning that there is no '$<$' relation.
\end{theorem}
Now, we work on several different problems:
\begin{enumerate}
  \item Express the following complex numbers for $a+ bi$.
  \begin{enumerate}
    \item $(a+3i)+(4-i)$
    \[(4 + a)-2i\]
    \item $(2+3i)(4-i)$
    \begin{gather*}
      8 + 3 - 2i + 12 i\\
      11 + 10i
    \end{gather*}
    \item $\displaystyle\frac{2+3i}{4-i}$
    \begin{align*}
      \frac{(2+3i)}{4-i}&= \frac{(2+3i)(4+i)}{(4-i)(4+i)}\\
      &= \left(\frac{5 + 14i}{15}\right)
    \end{align*}
  \end{enumerate}
  \item Find all solutions to $x^2 = 3 + 4i$. 
\end{enumerate}
\begin{definition}{Complex Conjugate}{}\gline{}
  The \ul{complex conjugate} of $z = a+bi$ is the number $\bar z = a - bi$. 
\end{definition}

\begin{definition}{Real \& Imaginary Part}{}\gline{}
  The \ul{real part} of $z = a+bi$ is $\real{z} = a$. \\
  The \ul{imaginary part} of $z = a+bi$ is $\im{z} = b$. \gline{}
  \textit{Note: }$\real{z},\im{z}\in \R$. 
\end{definition}

\begin{enumerate}
  \item [2.] Find all solutions to $z^2 = 3+ 4i$.
  \begin{problem}
    We want to use the formula $z = (a+bi)$. Thus, we get that:
    \begin{align*}
      z^2 = a^2 + 2abi - b^2
    \end{align*}
    Then, we find that:
    \begin{align*}
      3&= a^2 - b^2\\
      4 &= 2(ab) \quad\rightarrow \quad 2 &= (ab)
    \end{align*}
    So we get the solutions $(a,b) = (2, 1), (-2, -1)$.
  \end{problem}
  \item [3.] Find all solutions to $x^3 = 1$. 
  \begin{problem}
    We represent $z$ as $z = (a+bi)$. Thus, we find that:
    \begin{align*}
      z^3 &= a^3 + 3 a^2 b i + 3 a b^2 i^2 + b^3 i^3\\
      &= a^3 - 3ab + 3a^2bi - b^3 i
    \end{align*}
    And we get the following system of equations:
    \begin{align*}
      a^3 - 3ab &= 1\\
      3a^2b - b^3 &= 0
    \end{align*}
    For $a, b\ne 0$, we find that $3a^2 = b^2$. WLOG $a = b$ or $a = -b$. 
  \end{problem}
\end{enumerate}
\end{document}